\documentclass[11pt]{article}
\usepackage{manual_docs_style}
\externaldocument{build/terminology_shared_utilities}[terminology_shared_utilities.pdf]
\externaldocument{build/environment_profiles}[environment_profiles.pdf]
\externaldocument{build/simulink_generated_metadata_schema}[simulink_generated_metadata_schema.pdf]
\externaldocument{build/simulation_stage}[simulation_stage.pdf]

\begin{document}

\docTitle{Candidate Run Generation}
\phantomsection\label{docstart:candidate_generation}
This stage samples candidate run recipes according to the configured generation policies in \hyperref[sec:generation-policy-schema]{\code{generation_policies/<policy_id>.json}}.
Samples are generated deterministically from a seed. 
That is, given the same policy, simulator config, and seed, the same recipes will be produced.
The output is a resolved run graph that defines the runnable cases later materialized by \hyperref[docstart:simulation_stage]{Stage: Simulate}.

\section*{Dependencies}

For each \termSimulator{} under \code{models/simulink/<simulator_id>/}, the stage reads:

\begin{itemize}
\item \code{experiment_config.json}: hard environment contract, including exposed variables, legal intervals, simulation horizon, observed signals, and detectability configuration.
\item \code{generated/metadata.json}: generated Simulink metadata, including exposed workspace variables and intervention block mappings.
\end{itemize}

The stage writes a resolved plan under:

\begin{DocCode}
models/simulink/<simulator_id>/plans/<policy_id>/
\end{DocCode}

with these artifacts:

\begin{DocCode}
run_nodes.jsonl
run_edges.jsonl
generation_report.json
\end{DocCode}

\section*{Manual Cases}

\code{manual_cases.jsonl} is an optional high-level authoring format for curated cases.
Users specify baseline parameters and optional interventions; the candidate-generation compiler validates each case, expands it into runnable nodes, computes stable IDs and hashes, and writes the corresponding edges.
Users should not manually author \code{run_nodes.jsonl} or \code{run_edges.jsonl}.

Manual-only graph compilation can be run with:

\begin{DocCode}
python workflows/generate_runs/compile_run_graph.py \
  --model BallDrop \
  --policy manual_debug_v1 \
  --manual-only \
  --manual-cases manual_cases.jsonl \
  --overwrite
\end{DocCode}

Each line of \code{manual_cases.jsonl} is one JSON object of the following form:

\begin{DocCode}
{
“case_id”: “gravity_late_large_increase”,
“baseline_parameters”: {
“mass”: 4.0,
“drag_coeff”: 0.05,
“restitution”: 0.8,
“gravity”: 9.81,
“initial_height”: 3.0,
“initial_velocity”: -2.0
},
“interventions”: [
{
“parameter”: “gravity”,
“value”: 13.2,
“time”: 4.0
}
],
“notes”: “Curated case with late gravity increase.”
}
\end{DocCode}

During compilation, each manual case defines one baseline family.
For every intervention in \code{interventions}, the compiler creates the matched \code{intervention} node and \code{time0_baseline} node, then writes the corresponding \code{baseline_to_intervention} and \code{intervention_to_time0_baseline} edges.

The stage writes a plan under:

\begin{DocCode}
models/simulink/<simulator_id>/plans/<policy_id>/
\end{DocCode}

with these artifacts:

\begin{center}
\begin{tabular}{p{0.30\linewidth}p{0.60\linewidth}}
\hline
\textbf{Artifact} & \textbf{Purpose} \\
\hline
\code{run_nodes.jsonl} & Flat list of runnable simulation recipes. Each line defines one baseline, intervention, or time-zero/static-change run. \\
\code{run_edges.jsonl} & Graph relationships among runs, such as baseline-to-intervention and intervention-to-time-zero links. \\
\code{generation_report.json} & Generation summary, including requested counts, generated counts, skipped parameter-direction requests, and coverage statistics. \\
\hline
\end{tabular}
\end{center}

\section*{Execution Flow}

The candidate-generation stage performs the following steps:

\begin{enumerate}
\item Load \code{experiment_config.json}, \code{generation_policies/<policy_id>.json} validate internal consistency among files.
\begin{itemize}
  \item that all exposed variables are present in \code{generated/metadata.json}.
  \item \code{generation_policies/<policy_id>.json} has the sampling varables correct
\end{itemize}
\item Resolve the random seed and initialize deterministic random-number generators.
\item Sample baseline families according to \code{baseline_sampler}.
\item For each baseline family, sample intervention children according to \code{intervention_sampler}.
\item For each intervention child, create a matched time-zero/static-change recipe.
\item Canonicalize every recipe and compute \code{recipe_hash} and \code{run_id}.
\item Write flat runnable nodes to \code{run_nodes.jsonl}.
\item Write baseline/intervention/time-zero relationships to \code{run_edges.jsonl}.
\item Validate the resolved graph and write \code{validation_report.json}.
\end{enumerate}

The paper describes the third kind as the static-change probe.
The implementation may keep the name \code{time0_baseline} for backward compatibility.

\section*{\texorpdfstring{\titlecode{run_nodes.jsonl}}{run\_nodes.jsonl}}
\phantomsection\label{sec:candidate-generation-run-nodes}
Each line of \code{run_nodes.jsonl} is a complete runnable simulation run.
The resolved plan gives every run a stable \code{run_id}, gives every recipe a \code{recipe_hash}, and records baseline/intervention/time-zero relationships in \code{run_edges.jsonl}.

\subsection*{Minimal Shape}

\begin{DocCode}
{
  "run_id": "run_7fd9daa0055fb3f412ef5e0ab8d4c5ba",
  "kind": "intervention",
  "family_id": "run_7fd9daa0055fb3f412ef5e0ab8d4c5ba",
  "recipe_hash": "7fd9daa0055fb3f412ef5e0ab8d4c5ba",
  "recipe": {
    "model": "BallDrop",
    "baseline_parameters": {
      "mass": 4.0,
      "drag_coeff": 0.05,
      "restitution": 0.8,
      "gravity": 9.81,
      "initial_height": 3.0,
      "initial_velocity": -2.0
    },
    "intervention": {
      "parameter": "gravity",
      "value": 13.2,
      "time": 4.0
    }
  },
  "metadata": {
    "source": "programmatic",
    "policy_id": "production_v1",
    "seed": 456,
    "generator": "compile_run_graph.py"
  }
}
\end{DocCode}

\section*{Node Field Definitions}

\subsection*{\texorpdfstring{\titlecode{run_id}}{run\_id}}

\code{run_id} is the stable identifier for one runnable recipe.

Recommended form:

\begin{DocCode}
run_<32-hex-recipe-hash>
\end{DocCode}

The \code{run_id} must be unique within a plan.
Preferably, it should be globally stable across plans: if two policies produce the exact same canonical recipe, they should produce the same \code{run_id}.

\subsection*{\texorpdfstring{\titlecode{kind}}{kind}}

Allowed values are \code{baseline}, \code{intervention}, and \code{time0_baseline}.
Baseline nodes have no intervention; intervention nodes apply the changed value at \code{intervention.time}; time-zero/static-change nodes apply the changed value from time zero.

\subsection*{\texorpdfstring{\titlecode{family_id}}{family\_id}}

\code{family_id} groups one baseline run with its derived intervention and time-zero/static-change runs.
All nodes generated from the same baseline recipe should share the same \code{family_id}.

\subsection*{\texorpdfstring{\titlecode{recipe_hash}}{recipe\_hash}}

\code{recipe_hash} is the canonical hash of the \code{recipe} object only.
It must not include \code{metadata}.
Hashing should use normalized JSON values, canonical JSON with sorted keys and compact separators, SHA-256, and the first 32 hexadecimal characters of the digest.

\subsection*{\texorpdfstring{\titlecode{metadata}}{metadata}}

\code{metadata} records provenance for the node.
It is not part of the recipe hash.

Fields include:

\begin{itemize}
\item \code{source}: typically \code{programmatic} or \code{manual}.
\item \code{policy_id}: policy that produced the node.
\item \code{seed}: generation seed, when applicable.
\item \code{generator}: generator script name.
\end{itemize}

\section*{\texorpdfstring{\titlecode{recipe}}{recipe} Schema}

Each node's \code{recipe} defines the exact simulator input.
It contains:

\begin{itemize}
\item \code{model}: \termSimulator identifier.
\item \code{baseline_parameters}: all parameters and initial-state variables required to initialize the run.
\item \code{intervention}: changed parameter, post-intervention value, and intervention time. The top-level \code{kind} determines how these fields are interpreted.
\end{itemize}

\subsection*{Baseline Recipe}\label{sec:candidate-generation-baseline-recipe}

A baseline run has \code{kind = "baseline"} and no intervention.

\begin{DocCode}
{
  "model": "BallDrop",
  "baseline_parameters": {
    "mass": 4.0,
    "drag_coeff": 0.05,
    "restitution": 0.8,
    "gravity": 9.81,
    "initial_height": 3.0,
    "initial_velocity": -2.0
  },
  "intervention": {
    "parameter": null,
    "value": null,
    "time": null
  }
}
\end{DocCode}

\subsection*{Intervention Recipe}

An intervention run has \code{kind = "intervention"}.
It applies the changed value at \code{intervention.time}.

\begin{DocCode}
{
  "model": "BallDrop",
  "baseline_parameters": {
    "mass": 4.0,
    "drag_coeff": 0.05,
    "restitution": 0.8,
    "gravity": 9.81,
    "initial_height": 3.0,
    "initial_velocity": -2.0
  },
  "intervention": {
    "parameter": "gravity",
    "value": 13.2,
    "time": 4.0
  }
}
\end{DocCode}

\subsection*{Time-Zero / Static-Change Recipe}\label{sec:candidate-generation-time0-baseline-recipe}

A time-zero/static-change run has \code{kind = "time0_baseline"}.
It uses the same changed parameter and post-intervention value as the matched intervention run, but the value is active from time zero.

\begin{DocCode}
{
  "model": "BallDrop",
  "baseline_parameters": {
    "mass": 4.0,
    "drag_coeff": 0.05,
    "restitution": 0.8,
    "gravity": 9.81,
    "initial_height": 3.0,
    "initial_velocity": -2.0
  },
  "intervention": {
    "parameter": "gravity",
    "value": 13.2,
    "time": 0.0
  }
}
\end{DocCode}

\section*{\texorpdfstring{\titlecode{baseline_parameters}}{baseline\_parameters}}

\code{baseline_parameters} contains all simulator variables required to initialize the run.
It must include every variable listed in:

\begin{DocCode}
experiment_config.json::exposed_variables.parameters
experiment_config.json::exposed_variables.initial_state
\end{DocCode}

All values must be finite JSON numbers.
Each value must lie inside the corresponding \code{allowed_intervals} from \code{experiment_config.json}.

Variables under \code{parameters} may be intervened on.
Variables under \code{initial_state} are baseline-only variables and must not change after simulation start.

\section*{\texorpdfstring{\titlecode{intervention}}{intervention}}

\code{intervention} defines whether and how the baseline recipe is changed.

Fields:

\begin{itemize}
\item \code{parameter}: changed root-cause parameter, or \code{null} for baseline runs.
\item \code{value}: post-intervention value, or \code{null} for baseline runs.
\item \code{time}: intervention time, \code{0.0} for time-zero/static-change runs, or \code{null} for baseline runs.
\end{itemize}

Validation depends on top-level \code{kind}:

\begin{itemize}
\item If \code{kind = "baseline"}, \code{intervention.parameter}, \code{intervention.value}, and \code{intervention.time} must be \code{null}.
\item If \code{kind = "intervention"}, \code{intervention.parameter} and \code{intervention.value} must be non-\code{null}, and \code{intervention.time} must satisfy:
\end{itemize}

\begin{DocCode}
0 < intervention.time < experiment_config.json::end_time_input_s
\end{DocCode}

\begin{itemize}
\item If \code{kind = "time0_baseline"}, \code{intervention.parameter} and \code{intervention.value} must match the paired intervention, and \code{intervention.time} must be \code{0.0}.
\end{itemize}

Stage: Simulate reads top-level \code{kind} together with \code{recipe.intervention}.
It must not read \code{recipe.intervention.mode}.

\section*{\texorpdfstring{\titlecode{run_edges.jsonl}}{run\_edges.jsonl}}
\phantomsection\label{sec:candidate-generation-run-edges}

\code{run_edges.jsonl} is a JSON Lines file.
Each line is one relationship between run nodes.
Edges are used to recover the baseline/intervention/time-zero structure without nesting nodes.

\subsection*{Edge Shape}

\begin{DocCode}
{
  "edge_type": "baseline_to_intervention",
  "family_id": "run_7fd9daa0055fb3f412ef5e0ab8d4c5ba",
  "source_run_id": "run_baseline_...",
  "target_run_id": "run_intervention_...",
  "metadata": {
    "parameter": "gravity",
    "direction": "increase",
    "intervention_time": 4.0
  }
}
\end{DocCode}

Allowed \code{edge_type} values are:

\begin{itemize}
\item \code{baseline_to_intervention}
\item \code{intervention_to_time0_baseline}
\end{itemize}

Family membership is primarily represented by \code{family_id}.
A separate family-membership edge is not required.

\section*{\texorpdfstring{\titlecode{generation_report.json}}{generation\_report.json}}

\code{generation_report.json} records what the candidate-generation stage attempted and produced.

Recommended fields include:

\begin{itemize}
\item \code{timestamp}
\item \code{model}: \termSimulator identifier
\item \code{policy_id}
\item \code{seed}
\item \code{generator}
\item \code{generator_version}
\item requested family count
\item generated family count
\item baseline node count
\item intervention node count
\item time-zero/static-change node count
\item skipped parameter-direction requests
\item per-parameter and per-direction coverage summaries
\item resolved intervention value-sampler configuration, including \code{distance_space}
\item policy acceptance summary
\end{itemize}


\section*{Core Validation Rules}

The following rules apply to every resolved run graph with reference to \code{  run_nodes.jsonl}, \code{  run_edges.jsonl} and \code{validation_report.json}

\begin{itemize}
\item Every baseline recipe must include all exposed parameters and initial-state variables.
\item Every \code{kind = "baseline"} node must have \code{intervention.parameter}, \code{intervention.value}, and \code{intervention.time} set to \code{null}.
\item Every \code{kind = "intervention"} node must change exactly one variable listed under \code{experiment_config.json::exposed_variables.parameters}.
\item No intervention may change a variable listed only under \code{experiment_config.json::exposed_variables.initial_state}.
\item Every intervention value must respect the changed parameter's \code{allowed_intervals}.
\item Every edge must point to existing run IDs.
\item Every \code{baseline_to_intervention} edge must connect a baseline node to an intervention node.
\item Every \code{intervention_to_time0_baseline} edge must connect an intervention node to a time-zero/static-change node.
\end{itemize}

\section*{Additional Candidate Validation Rules}

The following additional checks apply to candidate runs intended for benchmark selection:

\begin{itemize}
 \item Every intervention candidate should have both a matched no-intervention baseline run and a matched time-zero/static-change run.
\end{itemize}


\end{document}
